% =================================================================
%  CSE 815 -- Machine Learning 
%  Module 1, Week 1, Part 2   
%  Visualising the Cost Function and Gradient Descent
% =================================================================
\documentclass[aspectratio=169,10pt]{beamer}
\usepackage{cse815}

\title[CSE 815 -- Week 1, Part 2]{Training a Model with Gradient Descent}
\subtitle{Visualising $J(w,b)$, the Descent Algorithm, and the Learning Rate}
\author{Rokan Uddin Faruqui\\ \small Professor, Department of CSE}
\institute{\small University of Chittagong \\ \texttt{rufaruqui@cu.ac.bd}}
\date[Week 1 / P 2]{ Week 1 $\cdot$ Part 2 }

\begin{document}

\begin{frame}[plain,noframenumbering]
  \titlepage
\end{frame}

% ---------------------------------------------------------------
\begin{frame}{Where we left off, and where we are going}
  \begin{columns}[T,onlytextwidth]
    \begin{column}{0.5\textwidth}
      \begin{block}{Part 1 gave us}
        \[ \fwb(x)=wx+b \]
        \[ J(w,b)=\frac{1}{2m}\sum_{i=1}^{m}\bigl(\fwb(\xii{i})-\yii{i}\bigr)^2 \]
        \[ \text{Goal:}\quad \min_{w,b} J(w,b) \]
      \end{block}
      \small We can \emph{state} the goal. We cannot yet \emph{reach} it: trying
      every $(w,b)$ is impossible --- the parameter space is continuous.
    \end{column}
    \begin{column}{0.46\textwidth}
      \begin{keyidea}[Today]
        \small
        \begin{enumerate}
          \item See $J$ as a surface and as contours 
          \item Gradient descent: the algorithm 
          \item Why the derivative points the right way 
          \item The learning rate $\alpha$ 
          \item Derive the gradients for linear regression 
          \item Batch gradient descent in action 
        \end{enumerate}
      \end{keyidea}
    \end{column}
  \end{columns}
\end{frame}

% ===============================================================
\section{Visualising the Cost Function}

\begin{frame}{With two parameters, $J$ is a surface}
  \begin{columns}[T,onlytextwidth]
    \begin{column}{0.55\textwidth}
      \centering
      \begin{tikzpicture}
        \begin{axis}[width=6.6cm,height=4.4cm,
          view={40}{30},
          xlabel={$w$},ylabel={$b$},zlabel={$J(w,b)$},
          tick label style={font=\tiny},label style={font=\scriptsize},
          colormap={cu}{color=(culight) color=(cuteal) color=(cublue)},
          samples=26,domain=-1.2:3.2,y domain=-1.6:3.6]
          \addplot3[surf,shader=interp,opacity=0.92]
            {1.6*(x-1)^2 + 0.9*(y-1)^2 + 0.55*(x-1)*(y-1) + 0.2};
        \end{axis}
      \end{tikzpicture}
    \end{column}
    \begin{column}{0.42\textwidth}
      \footnotesize
      \begin{itemize}
        \item Horizontal axes: the \emph{parameters} $w$ and $b$.
        \item Vertical axis: the cost $J$ --- how badly that line fits the
              \emph{fixed} training data.
        \item A \textbf{convex bowl}: one lowest point, no local minima, no
              plateaus.
        \item Each point on the surface is one candidate line $\fwb$.
      \end{itemize}
      \begin{keyidea}
        Training searches for the bottom of this bowl, blindfolded: we feel the
        local slope but never see the whole surface.
      \end{keyidea}
    \end{column}
  \end{columns}
\end{frame}

\begin{frame}{The same surface, seen from above: contour plots}
  \begin{columns}[T,onlytextwidth]
    \begin{column}{0.52\textwidth}
      \centering
      \begin{tikzpicture}
        \begin{axis}[width=6.4cm,height=5.0cm,xlabel={$w$},ylabel={$b$},
          tick label style={font=\scriptsize},label style={font=\scriptsize},
          view={0}{90}, domain=-1.2:3.2, y domain=-1.6:3.6,
          colormap={cu}{color=(cublue) color=(cuteal) color=(cuamber)}]
          \costcontour{}{0.3,0.8,1.6,2.8,4.4,6.4}{1.6*(x-1)^2 + 0.9*(y-1)^2 + 0.55*(x-1)*(y-1) + 0.2}
          \addplot3[only marks,mark=*,mark size=2.4pt,curose] coordinates {(1,1,0.2)};
          \node[font=\scriptsize,curose,anchor=west] at (axis cs:1.35,0.55,0) {minimum};
        \end{axis}
      \end{tikzpicture}
    \end{column}
    \begin{column}{0.44\textwidth}
      \small
      \begin{itemize}
        \item Each closed curve joins all $(w,b)$ pairs having the \emph{same}
              cost --- exactly like elevation lines on a topographic map of
              Chittagong Hill Tracts.
        \item The centre of the innermost ellipse is the minimum.
        \item Contours \emph{close together} $\Rightarrow$ steep slope;
              \emph{far apart} $\Rightarrow$ nearly flat.
        \item Elongated ellipses signal that the two parameters have very
              different scales --- this motivates feature scaling in Week 2.
      \end{itemize}
    \end{column}
  \end{columns}
\end{frame}

\begin{frame}{Reading the contour map: three candidate lines}
  \begin{columns}[T,onlytextwidth]
    \begin{column}{0.5\textwidth}
      \centering {\scriptsize Data and three fits}\par
      \begin{tikzpicture}
        \begin{axis}[width=6.0cm,height=4.4cm,xlabel={$x$},ylabel={$y$},
          xmin=0,xmax=9,ymin=0,ymax=10,grid=major,grid style={culight},
          tick label style={font=\scriptsize},label style={font=\scriptsize},
          legend style={font=\tiny,at={(0.02,0.98)},anchor=north west,draw=none,fill=white}]
          \addplot[only marks,mark=*,mark size=1.8pt,cublue,forget plot]
            coordinates {(1,2)(2,3.4)(3,3.8)(4,5.6)(5,5.4)(6,7.2)(7,7.4)(8,8.8)};
          \addplot[cuteal,very thick,domain=0.5:8.5]{0.95*x+1.1};   \addlegendentry{A: good}
          \addplot[cuamber,very thick,dashed,domain=0.5:8.5]{0.55*x+2.6}; \addlegendentry{B: fair}
          \addplot[curose,very thick,dotted,domain=0.5:8.5]{0.1*x+7}; \addlegendentry{C: poor}
        \end{axis}
      \end{tikzpicture}
    \end{column}
    \begin{column}{0.48\textwidth}
      \centering {\scriptsize Their positions in parameter space}\par
      \begin{tikzpicture}
        \begin{axis}[width=6.0cm,height=4.4cm,xlabel={$w$},ylabel={$b$},
          tick label style={font=\scriptsize},label style={font=\scriptsize},
          view={0}{90}, domain=-0.2:1.6, y domain=-1:8,
          colormap={cu}{color=(cublue) color=(cuteal) color=(cuamber)}]
          \costcontour{}{0.15,0.6,1.6,3.4,6.5,11}{14*(x-0.95)^2 + 0.42*(y-1.1)^2 + 3.2*(x-0.95)*(y-1.1) + 0.05}
          \addplot3[only marks,mark=*,mark size=2.2pt,cuteal] coordinates {(0.95,1.1,0)};
          \addplot3[only marks,mark=square*,mark size=2.2pt,cuamber] coordinates {(0.55,2.6,0)};
          \addplot3[only marks,mark=triangle*,mark size=2.8pt,curose] coordinates {(0.1,7,0)};
          \node[font=\tiny,cuteal,anchor=west]  at (axis cs:1.02,1.1,0) {A};
          \node[font=\tiny,cuamber,anchor=west] at (axis cs:0.62,2.6,0) {B};
          \node[font=\tiny,curose,anchor=west]  at (axis cs:0.17,7,0) {C};
        \end{axis}
      \end{tikzpicture}
    \end{column}
  \end{columns}
  \vspace{1mm}
  \small
  \centering
  A sits at the centre (low cost); C sits on an outer contour (high cost).
  \textbf{Better line $\Leftrightarrow$ inner contour.}
\end{frame}

% ===============================================================
\section{Gradient Descent}

\begin{frame}{The idea: walk downhill}
  \begin{columns}[T,onlytextwidth]
    \begin{column}{0.52\textwidth}
      \small
      You are standing on a hillside in dense fog. You want the valley floor.
      A reasonable procedure:
      \begin{enumerate}
        \item Feel the ground around your feet.
        \item Identify the direction of \emph{steepest descent}.
        \item Take a small step that way.
        \item Repeat until the ground is flat.
      \end{enumerate}
      \vspace{2mm}
      That procedure is \textbf{gradient descent}. The ``feel the ground'' step
      is a derivative; the ``small step'' is scaled by the learning rate.
    \end{column}
    \begin{column}{0.46\textwidth}
      \centering
      \begin{tikzpicture}
        \begin{axis}[width=6.0cm,height=4.6cm,
          xlabel={$w$},ylabel={$J(w)$},
          xmin=-2.4,xmax=3.4,ymin=-0.4,ymax=6.4,
          grid=major,grid style={culight},
          tick label style={font=\scriptsize},label style={font=\scriptsize}]
          \addplot[cublue,very thick,domain=-2.4:3.4,samples=90]{0.6*(x-0.5)^2+0.4};
          \addplot[only marks,mark=*,mark size=2.2pt,curose]
            coordinates {(-2.0,4.15)(-1.0,1.75)(-0.4,0.886)(-0.04,0.578)(0.18,0.461)};
          \draw[-{Stealth},cuamber,thick] (axis cs:-2.0,4.15)--(axis cs:-1.0,1.75);
          \draw[-{Stealth},cuamber,thick] (axis cs:-1.0,1.75)--(axis cs:-0.4,0.886);
          \draw[-{Stealth},cuamber,thick] (axis cs:-0.4,0.886)--(axis cs:-0.04,0.578);
          \node[font=\scriptsize,cugray] at (axis cs:2.2,1.0) {steps shrink};
        \end{axis}
      \end{tikzpicture}
    \end{column}
  \end{columns}
\end{frame}

\begin{frame}{The gradient descent algorithm}
  \begin{defnbox}[Gradient descent for two parameters]
    Start from some initial $w,b$ (commonly $w=0,\,b=0$). Repeat until convergence:
    \[
      \left\{
      \begin{aligned}
        \text{tmp}_w &:= w - \alpha\,\frac{\partial}{\partial w}J(w,b)\\[2pt]
        \text{tmp}_b &:= b - \alpha\,\frac{\partial}{\partial b}J(w,b)\\[2pt]
        w &:= \text{tmp}_w,\qquad b := \text{tmp}_b
      \end{aligned}
      \right.
    \]
  \end{defnbox}
  \small
  \begin{itemize}
    \item $\alpha>0$ is the \textbf{learning rate}, typically $0.001$--$0.3$; it
          controls the step size.
    \item $:=$ denotes assignment, not a mathematical equality.
    \item The minus sign is what makes it \emph{descent}: we move \emph{against}
          the gradient, the direction of steepest increase.
  \end{itemize}
\end{frame}

\begin{frame}{Simultaneous update --- a correctness requirement}
  \begin{columns}[T,onlytextwidth]
    \begin{column}{0.48\textwidth}
      \begin{block}{Correct: simultaneous}
        \begin{align*}
          \text{tmp}_w &:= w-\alpha\,\partial_w J(w,b)\\
          \text{tmp}_b &:= b-\alpha\,\partial_b J(w,b)\\
          w &:= \text{tmp}_w\\
          b &:= \text{tmp}_b
        \end{align*}
        \footnotesize Both derivatives are evaluated at the \emph{same} point
        $(w,b)$ --- the old one.
      \end{block}
    \end{column}
    \begin{column}{0.48\textwidth}
      \begin{alertblock}{Incorrect: sequential}
        \begin{align*}
          w &:= w-\alpha\,\partial_w J(w,b)\\
          b &:= b-\alpha\,\partial_b J(\underbrace{w}_{\text{already changed!}},b)
        \end{align*}
        \footnotesize The second derivative is taken at a \emph{different} point.
        This is another algorithm; it sometimes still converges, which makes the
        bug hard to spot.
      \end{alertblock}
    \end{column}
  \end{columns}
  \begin{pitfall}
    \footnotesize In NumPy, \texttt{w -= alpha*dj\_dw} then
    \texttt{b -= alpha*dj\_db} is safe \emph{only if} both gradients were
    computed \emph{before} either assignment.
  \end{pitfall}
\end{frame}

\begin{frame}{Why the derivative points the right way}
  \framesubtitle{One parameter, $J(w)$, update $w := w - \alpha\,\frac{d}{dw}J(w)$}
  \begin{columns}[T,onlytextwidth]
    \begin{column}{0.5\textwidth}
      \centering
      \begin{tikzpicture}
        \begin{axis}[width=6.2cm,height=4.8cm,xlabel={$w$},ylabel={$J(w)$},
          xmin=-2.6,xmax=3.6,ymin=-0.3,ymax=5.6,
          grid=major,grid style={culight},
          tick label style={font=\scriptsize},label style={font=\scriptsize}]
          \addplot[cublue,very thick,domain=-2.6:3.6,samples=90]{0.5*(x-0.5)^2+0.3};
          % right point, positive slope
          \addplot[only marks,mark=*,mark size=2.2pt,curose] coordinates {(2.7,2.72)};
          \addplot[curose,thick,domain=1.9:3.5]{2.2*(x-2.7)+2.72};
          \draw[-{Stealth},curose,very thick] (axis cs:2.7,2.1)--(axis cs:1.9,2.1);
          % left point, negative slope
          \addplot[only marks,mark=*,mark size=2.2pt,cuteal] coordinates {(-1.7,2.72)};
          \addplot[cuteal,thick,domain=-2.5:-0.9]{-2.2*(x+1.7)+2.72};
          \draw[-{Stealth},cuteal,very thick] (axis cs:-1.7,2.1)--(axis cs:-0.9,2.1);
        \end{axis}
      \end{tikzpicture}
    \end{column}
    \begin{column}{0.46\textwidth}
      \small
      \textbf{Right of the minimum} ({\color{curose}rose}): the tangent rises, so
      $\frac{dJ}{dw}>0$. Then
      \[ w := w - \alpha\cdot(\text{positive}) \;\Rightarrow\; w\ \text{decreases}. \]
      We move \emph{left}, toward the minimum. \pause
      \vspace{2mm}

      \textbf{Left of the minimum} ({\color{cuteal}teal}): the tangent falls, so
      $\frac{dJ}{dw}<0$. Then
      \[ w := w - \alpha\cdot(\text{negative}) \;\Rightarrow\; w\ \text{increases}. \]
      We move \emph{right}, toward the minimum.\pause
      \vspace{2mm}

      \textbf{At the minimum:} $\frac{dJ}{dw}=0$, so $w:=w-0$. Gradient descent
      leaves a minimum alone --- it is a fixed point.
    \end{column}
  \end{columns}
\end{frame}

% ===============================================================
\section{The Learning Rate}

\begin{frame}{Choosing $\alpha$: too small, too large, just right}
  \begin{columns}[T,onlytextwidth]
    \begin{column}{0.33\textwidth}
      \centering {\scriptsize $\alpha$ too \textbf{small}}\par
      \begin{tikzpicture}
        \begin{axis}[width=4.3cm,height=3.6cm,xtick=\empty,ytick=\empty,
          xmin=-2.6,xmax=2.6,ymin=0,ymax=4,axis lines=box]
          \addplot[cublue,thick,domain=-2.6:2.6,samples=70]{0.55*x^2+0.2};
          \addplot[only marks,mark=*,mark size=1.5pt,curose]
            coordinates {(-2.2,2.86)(-1.98,2.36)(-1.78,1.94)(-1.6,1.61)(-1.44,1.34)(-1.3,1.13)};
          \draw[-{Stealth},cuamber] (axis cs:-2.2,2.86)--(axis cs:-1.98,2.36);
          \draw[-{Stealth},cuamber] (axis cs:-1.98,2.36)--(axis cs:-1.78,1.94);
          \draw[-{Stealth},cuamber] (axis cs:-1.78,1.94)--(axis cs:-1.6,1.61);
        \end{axis}
      \end{tikzpicture}
      {\scriptsize Converges, but needs a huge number of iterations. Training is
      slow, not wrong.}
    \end{column}
    \begin{column}{0.33\textwidth}
      \centering {\scriptsize $\alpha$ too \textbf{large}}\par
      \begin{tikzpicture}
        \begin{axis}[width=4.3cm,height=3.6cm,xtick=\empty,ytick=\empty,
          xmin=-3.2,xmax=3.2,ymin=0,ymax=6,axis lines=box]
          \addplot[cublue,thick,domain=-3.2:3.2,samples=70]{0.55*x^2+0.2};
          \addplot[only marks,mark=*,mark size=1.5pt,curose]
            coordinates {(-0.8,0.55)(1.2,1.0)(-1.9,2.19)(2.9,4.83)};
          \draw[-{Stealth},curose] (axis cs:-0.8,0.55)--(axis cs:1.2,1.0);
          \draw[-{Stealth},curose] (axis cs:1.2,1.0)--(axis cs:-1.9,2.19);
          \draw[-{Stealth},curose] (axis cs:-1.9,2.19)--(axis cs:2.9,4.83);
        \end{axis}
      \end{tikzpicture}
      {\scriptsize Overshoots the minimum each step; $J$ \emph{increases}.
      Diverges.}
    \end{column}
    \begin{column}{0.33\textwidth}
      \centering {\scriptsize $\alpha$ \textbf{well chosen}}\par
      \begin{tikzpicture}
        \begin{axis}[width=4.3cm,height=3.6cm,xtick=\empty,ytick=\empty,
          xmin=-2.6,xmax=2.6,ymin=0,ymax=4,axis lines=box]
          \addplot[cublue,thick,domain=-2.6:2.6,samples=70]{0.55*x^2+0.2};
          \addplot[only marks,mark=*,mark size=1.5pt,cuteal]
            coordinates {(-2.2,2.86)(-1.1,0.87)(-0.55,0.37)(-0.27,0.24)(-0.13,0.21)};
          \draw[-{Stealth},cuteal] (axis cs:-2.2,2.86)--(axis cs:-1.1,0.87);
          \draw[-{Stealth},cuteal] (axis cs:-1.1,0.87)--(axis cs:-0.55,0.37);
          \draw[-{Stealth},cuteal] (axis cs:-0.55,0.37)--(axis cs:-0.27,0.24);
        \end{axis}
      \end{tikzpicture}
      {\scriptsize Rapid, monotone decrease in $J$.}
    \end{column}
  \end{columns}
  \vspace{2mm}
  \begin{keyidea}
    Try $\alpha \in \{0.001,\,0.003,\,0.01,\,0.03,\,0.1,\,0.3\}$ --- roughly
    $3\times$ steps. Plot $J$ against iteration number for each and pick the
    largest $\alpha$ that still decreases $J$ smoothly.
  \end{keyidea}
\end{frame}

\begin{frame}{Steps shrink automatically --- $\alpha$ need not}
  \begin{columns}[T,onlytextwidth]
    \begin{column}{0.5\textwidth}
      \centering
      \begin{tikzpicture}
        \begin{axis}[width=6.2cm,height=4.6cm,xlabel={$w$},ylabel={$J(w)$},
          xmin=-0.4,xmax=3.4,ymin=-0.2,ymax=4.4,
          grid=major,grid style={culight},
          tick label style={font=\scriptsize},label style={font=\scriptsize}]
          \addplot[cublue,very thick,domain=-0.4:3.4,samples=90]{0.9*(x-0.4)^2+0.15};
          \addplot[only marks,mark=*,mark size=2pt,curose]
            coordinates {(3.0,4.03)(1.96,2.28)(1.34,1.03)(0.96,0.46)(0.74,0.26)(0.6,0.19)};
          \draw[-{Stealth},cuamber,thick] (axis cs:3.0,4.03)--(axis cs:1.96,2.28);
          \draw[-{Stealth},cuamber,thick] (axis cs:1.96,2.28)--(axis cs:1.34,1.03);
          \draw[-{Stealth},cuamber,thick] (axis cs:1.34,1.03)--(axis cs:0.96,0.46);
          \draw[-{Stealth},cuamber,thick] (axis cs:0.96,0.46)--(axis cs:0.74,0.26);
        \end{axis}
      \end{tikzpicture}
    \end{column}
    \begin{column}{0.46\textwidth}
      \small
      As we approach the minimum the slope $\bigl|\partial J/\partial w\bigr|$
      gets smaller, so the step
      \[ \Delta w = -\alpha\,\frac{\partial J}{\partial w} \]
      gets smaller too --- \emph{with $\alpha$ held fixed}.
      \vspace{2mm}
      \begin{keyidea}
        Gradient descent decelerates near a minimum on its own. For a convex
        problem like linear regression, $\alpha$ need not be decayed by hand.
      \end{keyidea}
      \footnotesize For non-convex problems a fixed $\alpha$ can stall in a
      \emph{local} minimum --- Module 2 introduces Adam.
    \end{column}
  \end{columns}
\end{frame}

% ===============================================================
\section{Gradient Descent for Linear Regression}

\begin{frame}{Deriving the two partial derivatives}
  \small
  With $\fwb(x)=wx+b$ and $J(w,b)=\dfrac{1}{2m}\displaystyle\sum_{i=1}^{m}\bigl(\fwb(\xii{i})-\yii{i}\bigr)^2$:
  \vspace{1mm}
  \[
    \frac{\partial J}{\partial w}
    = \frac{1}{2m}\sum_{i=1}^{m} 2\bigl(w\xii{i}+b-\yii{i}\bigr)\cdot \frac{\partial}{\partial w}\bigl(w\xii{i}+b-\yii{i}\bigr)
    = \frac{1}{m}\sum_{i=1}^{m}\bigl(\fwb(\xii{i})-\yii{i}\bigr)\,\xii{i}
  \]
  \[
    \frac{\partial J}{\partial b}
    = \frac{1}{2m}\sum_{i=1}^{m} 2\bigl(w\xii{i}+b-\yii{i}\bigr)\cdot 1
    = \frac{1}{m}\sum_{i=1}^{m}\bigl(\fwb(\xii{i})-\yii{i}\bigr)
  \]
  \vspace{1mm}
  \begin{keyidea}
    This is where the $\tfrac{1}{2}$ in $J$ pays for itself: the $2$ from the
    chain rule cancels it exactly. The $\tfrac12$ changes the \emph{value} of $J$
    but not its \emph{minimiser}.
  \end{keyidea}
  \vspace{1mm}
  Note the structure: both gradients are averages of the residual
  $\bigl(\yhat^{(i)}-\yii{i}\bigr)$ --- weighted by $\xii{i}$ for $w$, unweighted for $b$.
\end{frame}

\begin{frame}{The complete algorithm}
  \begin{defnbox}[Batch gradient descent for univariate linear regression]
    Initialise $w:=0$, $b:=0$. Repeat until convergence:
    \[
      \begin{aligned}
        w &:= w - \alpha\,\frac{1}{m}\sum_{i=1}^{m}\bigl(\fwb(\xii{i})-\yii{i}\bigr)\xii{i}\\[3pt]
        b &:= b - \alpha\,\frac{1}{m}\sum_{i=1}^{m}\bigl(\fwb(\xii{i})-\yii{i}\bigr)
      \end{aligned}
      \qquad\text{(updated simultaneously)}
    \]
  \end{defnbox}
  \small
  \textbf{``Batch''} means each update uses \emph{all $m$} examples; mini-batch
  and stochastic variants use a subset per step.
  \begin{keyidea}[Why we are safe here]
    \footnotesize The squared-error cost of a linear model is \emph{convex} in
    $(w,b)$ --- a single bowl. With a suitable $\alpha$, gradient descent
    always reaches the \emph{global} minimum, whatever the initialisation.
  \end{keyidea}
\end{frame}

\begin{frame}{One iteration by hand}
  \small
  Data $(1,1),(2,2),(3,3)$, $m=3$. Start at $w=0,\,b=0$, take $\alpha=0.1$.
  \vspace{1mm}

  \textbf{Step 1 --- predictions and residuals.}
  $\hat y = (0,0,0)$, so residuals $\hat y^{(i)}-y^{(i)} = (-1,-2,-3)$.
  \vspace{1mm}

  \textbf{Step 2 --- gradients.}
  \[
    \frac{\partial J}{\partial w}=\frac{1}{3}\bigl[(-1)(1)+(-2)(2)+(-3)(3)\bigr]
      =\frac{-14}{3}\approx-4.667,
    \qquad
    \frac{\partial J}{\partial b}=\frac{-1-2-3}{3}=-2.
  \]
  \textbf{Step 3 --- simultaneous update.}
  \[
    w := 0 - 0.1(-4.667) = 0.467,\qquad b := 0 - 0.1(-2) = 0.2 .
  \]
  \vspace{1mm}
  \textbf{Check.} Cost fell from $J(0,0)=\tfrac{1}{6}(1+4+9)=2.33$ to
  $J(0.467,0.2)\approx 0.47$. Both parameters moved toward the true
  $(w,b)=(1,0)$; $b$ will be pulled back down in later iterations.
  \vspace{1mm}
  \begin{pitfall}
    Note that $b$ overshot \emph{upward} on iteration 1 even though its optimum is
    $0$. Individual parameters may move away from their optimum; only $J$ is
    guaranteed to decrease.
  \end{pitfall}
\end{frame}

\begin{frame}{Running gradient descent: the trajectory}
  \begin{columns}[T,onlytextwidth]
    \begin{column}{0.5\textwidth}
      \centering {\scriptsize Path over the contours of $J(w,b)$}\par
      \begin{tikzpicture}
        \begin{axis}[width=6.1cm,height=4.7cm,xlabel={$w$},ylabel={$b$},
          tick label style={font=\scriptsize},label style={font=\scriptsize},
          view={0}{90}, domain=-0.4:1.8, y domain=-1.2:4.4,
          colormap={cu}{color=(cublue) color=(cuteal) color=(cuamber)}]
          \costcontour{}{0.1,0.45,1.2,2.6,5,9}{9*(x-1)^2 + 0.5*(y-0.2)^2 + 2.4*(x-1)*(y-0.2) + 0.02}
          \addplot3[curose,thick,mark=*,mark size=1.5pt] coordinates {
            (0,3.6,0) (0.62,2.5,0) (0.86,1.6,0) (0.95,0.95,0) (0.99,0.5,0) (1.0,0.28,0)};
          \addplot3[only marks,mark=star,mark size=3.4pt,cublue] coordinates {(1,0.2,0)};
          \node[font=\tiny,cugray,anchor=west] at (axis cs:0.05,3.95,0) {start};
        \end{axis}
      \end{tikzpicture}
    \end{column}
    \begin{column}{0.48\textwidth}
      \centering {\scriptsize The learning curve: $J$ vs.\ iteration}\par
      \begin{tikzpicture}
        \begin{axis}[width=6.1cm,height=4.7cm,xlabel={iteration},ylabel={$J(w,b)$},
          xmin=0,xmax=60,ymin=0,ymax=2.6,grid=major,grid style={culight},
          tick label style={font=\scriptsize},label style={font=\scriptsize}]
          \addplot[cublue,very thick,domain=0:60,samples=90]{2.4*exp(-0.12*x)+0.02};
          \draw[dashed,cugray] (axis cs:0,0.05)--(axis cs:60,0.05);
          \node[font=\scriptsize,cugray,anchor=south east] at (axis cs:58,0.12) {plateau $\Rightarrow$ converged};
        \end{axis}
      \end{tikzpicture}
    \end{column}
  \end{columns}
  \begin{keyidea}[Always plot the learning curve]
    \footnotesize A correct gradient descent decreases $J$ on \emph{every}
    iteration. If $J$ ever rises, $\alpha$ is too large or the gradient code
    has a bug.
  \end{keyidea}
\end{frame}

\begin{frame}{Declaring convergence}
  \small
  \begin{columns}[T,onlytextwidth]
    \begin{column}{0.5\textwidth}
      \begin{block}{Automatic test}
        Stop when $J$ decreases by less than $\varepsilon$ (e.g.\ $10^{-3}$) in one
        iteration.\\[3pt]
        \emph{Drawback:} the right $\varepsilon$ varies wildly across problems and
        is hard to guess in advance.
      \end{block}
      \begin{block}{Visual test --- preferred}
        Plot $J$ against iteration count and look for the curve flattening.
        You also see divergence, oscillation and a badly chosen $\alpha$
        immediately.
      \end{block}
    \end{column}
    \begin{column}{0.48\textwidth}
      \begin{block}{Debugging checklist}
        \begin{itemize}
          \item $J$ increasing $\Rightarrow$ reduce $\alpha$ by $3$--$10\times$.
          \item $J$ oscillating $\Rightarrow$ $\alpha$ slightly too large.
          \item $J$ barely moving $\Rightarrow$ $\alpha$ too small, or features on
                wildly different scales (Week 2).
          \item Sanity check: with a very small $\alpha$, $J$ \emph{must} decrease
                every iteration. If it does not, the bug is in the gradient, not
                in $\alpha$.
        \end{itemize}
      \end{block}
    \end{column}
  \end{columns}
\end{frame}

% ===============================================================
\section{Wrap-up and Assessment}

\begin{frame}{Summary of Week 1}
  \small
  \begin{columns}[T,onlytextwidth]
    \begin{column}{0.5\textwidth}
      \begin{block}{The three-part recipe}
        \begin{enumerate}
          \item \textbf{Model:} $\fwb(x)=wx+b$
          \item \textbf{Cost:} $J(w,b)=\frac{1}{2m}\sum_i(\fwb(\xii{i})-\yii{i})^2$
          \item \textbf{Optimiser:} gradient descent
            \[ \theta := \theta - \alpha\,\partial_\theta J \]
        \end{enumerate}
        Every algorithm in this course --- logistic regression, neural networks,
        recommenders --- is a different choice of these same three parts.
      \end{block}
    \end{column}
    \begin{column}{0.48\textwidth}
      \begin{block}{Things to remember}
        \begin{itemize}
          \item Update $w$ and $b$ \emph{simultaneously}.
          \item The derivative's \emph{sign} steers; $\alpha$ sets the \emph{size}.
          \item Squared-error + linear model $\Rightarrow$ convex $\Rightarrow$
                global minimum guaranteed.
          \item ``Batch'' = all $m$ examples per step.
          \item Plot $J$ vs.\ iteration, always.
        \end{itemize}
      \end{block}
      \begin{block}{Next week}
        \small Multiple features, vectorisation, feature scaling, and choosing
        $\alpha$ systematically.
      \end{block}
    \end{column}
  \end{columns}
\end{frame}

\begin{frame}{Descriptive questions}
  \small
  \begin{enumerate}
    \item Explain what a contour plot of $J(w,b)$ represents. What does the
          spacing between adjacent contours tell you about the gradient there?
          \hfill\clo{5}
    \item Write the gradient descent update rule for $(w,b)$ and explain why the
          updates must be simultaneous. Show, with a small example, how a
          sequential update yields a different trajectory. \hfill\clo{6}
    \item Using the sign of $dJ/dw$, argue that the update
          $w := w-\alpha\,dJ/dw$ moves $w$ toward the minimum whether the current
          $w$ lies to the left or to the right of it. \hfill\clo{5}
    \item Derive $\partial J/\partial w$ and $\partial J/\partial b$ for
          univariate linear regression with the squared-error cost. State clearly
          where the chain rule is applied. \hfill\clo{8}
  \end{enumerate}
\end{frame}

\begin{frame}{Descriptive questions (continued)}
  \small
  \begin{enumerate}\setcounter{enumi}{4}
    \item Describe the behaviour of gradient descent when $\alpha$ is (a) too
          small, (b) too large. Sketch the corresponding curve of $J$ against
          iteration number for each case. \hfill\clo{5}
    \item Explain why gradient descent takes smaller steps near the minimum even
          though $\alpha$ is held constant. \hfill\clo{4}
    \item For the data $(1,1),(2,2),(3,3)$, model $\fwb(x)=wx+b$, initial
          $w=b=0$ and $\alpha=0.1$, carry out \emph{two} iterations of gradient
          descent by hand, reporting $w$, $b$ and $J$ after each. \hfill\clo{8}
    \item Why is the global optimum guaranteed for linear regression with squared
          error, but not for a neural network? \hfill\clo{4}
  \end{enumerate}
\end{frame}

\begin{frame}{Design questions}
  \small
  \begin{enumerate}
    \item \textbf{Design a learning-rate schedule.} You must train the same model
          on 200 different districts' data overnight, unattended --- no human can
          watch the learning curves. Design an automated procedure that selects
          $\alpha$ per district and detects divergence. Specify the decision rule,
          the quantities you would log, and the failure case your design does not
          handle. \hfill\clo{12}
    \item \textbf{Design a stopping criterion.} Propose a convergence test that is
          scale-invariant --- one that behaves the same whether prices are
          recorded in taka or in lakh taka. Justify why your criterion has this
          property, and contrast it with the naive test
          ``stop when $\Delta J < 10^{-3}$''. \hfill\clo{10}
    \item \textbf{Design a test suite.} A classmate's \texttt{compute\_gradient}
          function returns plausible numbers but the model never fits well. Design
          three unit tests that would localise the bug --- for example, tests
          exploiting a data set whose optimal parameters you know analytically.
          State what each test's expected output is and which bug it catches.
          \hfill\clo{12}
  \end{enumerate}
\end{frame}

\begin{frame}{Design questions (continued)}
  \small
  \begin{enumerate}\setcounter{enumi}{3}
    \item \textbf{Redesign the optimiser.} Suppose $\partial J/\partial w$ is
          expensive because $m = 10^{8}$. Propose a modification to batch gradient
          descent that reduces per-iteration cost, and analyse the trade-off it
          introduces in the convergence behaviour. \hfill\clo{10}
    \item \textbf{Diagnose from a plot.} You are shown a learning curve in which
          $J$ falls sharply for $50$ iterations, then rises steadily. List the two
          most likely causes and design an experiment that distinguishes between
          them. \hfill\clo{8}
  \end{enumerate}
\end{frame}

\begin{frame}{Reading}
  \begin{block}{Primary text}
    G.~James, D.~Witten, T.~Hastie and R.~Tibshirani,
    \emph{An Introduction to Statistical Learning}, 2nd ed., Springer, 2021.\\[2pt]
    \hfill{\small$\rightarrow$ \textbf{\S3.1.1--3.1.3} --- estimating the coefficients by
    least squares, and assessing the fit.}
  \end{block}
  \begin{block}{Supplementary}
    C.~M.~Bishop, \emph{Pattern Recognition and Machine Learning}, Springer, 2006.\\[2pt]
    \hfill{\small$\rightarrow$ \textbf{\S3.1.3} --- sequential (gradient-based) learning and the
    least-mean-squares update rule.}
  \end{block}
  \vspace{2mm}
  \begin{center}
    \small\color{cugray}
    Lab (CSE 816, Week 1): \emph{Gradient Descent} --- implement
    \texttt{compute\_cost} and \texttt{compute\_gradient}, then plot the learning
    curve.\\
    Practice quiz: \emph{Train the model with gradient descent}.
  \end{center}
\end{frame}

\end{document}
